\bibitem{b1} S. Wijethilaka and M. Liyanage, ``Survey on Network Slicing for Internet of Things Realization in 5G Networks,'' Apr. 01, 2021, \emph{Institute of Electrical and Electronics Engineers Inc.} doi: 10.1109/COMST.2021.3067807.

\bibitem{b2} V. Sharma, ``Performance Evaluation of Network Slicing in 5G Core Networks,'' \emph{International Journal of Multidisciplinary Research and Growth Evaluation}, vol. 3, no. 5, pp. 648--654, 2022, doi: 10.54660/.ijmrge.2022.3.5.648-654.

\bibitem{b3} E. Thomatos, A. Sgora, A. Tsipis, and P. Chatzimisios, ``AI Methods in Network Slice Life-Cycle Phases: A Survey,'' Oct. 01, 2025, \emph{Multidisciplinary Digital Publishing Institute (MDPI)}. doi: 10.3390/electronics14204053.

\bibitem{b4} W. He, H. Yao, H. Chang, and Y. Liu, ``A P4-Based Approach to Traffic Isolation and Bandwidth Management for 5G Network Slicing,'' \emph{Tsinghua Sci. Technol.}, vol. 30, no. 1, pp. 171--185, Sep. 2024, doi: 10.26599/tst.2024.9010020.

\bibitem{b5} Y. Safaei, P. Arefijamal, M. Siamaki, and B. Safaei, ``Priority-Aware SDN Orchestration for Surgical IoMT: A Joint Optimization of Hit Ratio and Latency With Dynamic Resource Reallocation,'' \emph{IEEE Access}, vol. 13, pp. 113787--113808, 2025, doi: 10.1109/ACCESS.2025.3583899.

\bibitem{b6} V. Onaji, E. Adediji, J. Njimgou, K. Ayano, and D. Olufemi, ``Deep Reinforcement Learning for Dynamic Network Slicing and Resource Orchestration in Software-Defined Critical Telecom Infrastructure,'' \emph{International Journal of Computer Applications Technology and Research}, Nov. 2025, doi: 10.7753/ijcatr1411.1006.

\bibitem{b7} P. Jangale~ , ``Impact of Network Slicing on Performance in 5G and Beyond,'' \emph{Journal of Artificial Intelligence, Machine Learning and Data Science}, vol. 1, no. 1, pp. 1653--1657, Nov. 2022, doi: 10.51219/jaimld/pratik-jangale/368.

\bibitem{b8} S. T. Arzo \emph{et al.}, ``MSN: A Playground Framework for Design and Evaluation of MicroServices-Based sdN Controller,'' \emph{Journal of Network and Systems Management}, vol. 30, no. 1, Jan. 2022, doi: 10.1007/s10922-021-09631-7.

\bibitem{b9} R. Pérez, M. Rivera, Y. Salgueiro, C. R. Baier, and P. Wheeler, ``Moving Microgrid Hierarchical Control to an SDN-Based Kubernetes Cluster: A Framework for Reliable and Flexible Energy Distribution,'' \emph{Sensors}, vol. 23, no. 7, Apr. 2023, doi: 10.3390/s23073395.

\bibitem{b10} V. Kumar Sharma, ``Cloud-Native 5G Deployments: Kubernetes and Microservices in Telco Networks,'' \emph{International Journal of Innovative Research in Engineering \& Multidisciplinary Physical Sciences}, vol. 10, no. 3, pp. 1--8, 2022, {[}Online{]}. Available: www.ijirmps.org

\bibitem{b11} A. Kamisetty, D. Narsina, M. Rodriguez, S. Kothapalli, J. Chandra, and S. Gummadi, ``Microservices vs. Monoliths: Comparative Analysis for Scalable Software Architecture Design,'' \emph{Engineering International}, vol. 11, no. 2, 2023.

\bibitem{b12} S. T. Arzo, D. Scotece, R. Bassoli, M. Devetsikiotis, L. Foschini, and F. H. P. Fitzek, ``Softwarized and containerized microservices-based network management analysis with MSN,'' \emph{Computer Networks}, vol. 254, Dec. 2024, doi: 10.1016/j.comnet.2024.110750.

\bibitem{b13} Y. Xu and Q. Yu, ``Deep reinforcement learning based computation offloading and resource allocation strategy for maritime internet of things,'' \emph{Computer Networks}, vol. 264, Jun. 2025, doi: 10.1016/j.comnet.2025.111221.

\bibitem{b14} B. M. Alrifai, L. Alsbatin, and F. Zawaideh, ``Optimizing 5G Resource Allocation with Reinforcement Learning: A Q-Learning and Deep Q-Network (DQN) Approach,'' \emph{Journal of Advances in Information Technology}, vol. 16, no. 11, pp. 1586--1594, 2025, doi: 10.12720/jait.16.11.1586-1594.

\bibitem{b15} M. Hosseinzadeh \emph{et al.}, ``SDN-Based NFV deployment for multi-objective resource allocation in edge computing: A deep reinforcement learning for iot workload scheduling,'' \emph{Sustainable Computing: Informatics and Systems}, vol. 48, Dec. 2025, doi: 10.1016/j.suscom.2025.101218.

\bibitem{b16} V. Sittakul, I. Ioannou, P. Nagaradjane, and V. Vassiliou, ``Intelligent congestion control in 5G URLLC Software-Defined Networks using adaptive resource management via Reinforced Dueling Deep Q-Networks,'' \emph{Journal of Network and Computer Applications}, vol. 242, Oct. 2025, doi: 10.1016/j.jnca.2025.104276.

\bibitem{b17} A. Iqbal, A. Nauman, T. Khurshaid, and S. B. Rhee, ``A Scalable Reinforcement Learning Framework for Ultra-Reliable Low-Latency Spectrum Management in Healthcare Internet of Things,'' \emph{Mathematics}, vol. 13, no. 18, Sep. 2025, doi: 10.3390/math13182941.

\bibitem{b18} M. Chahoud \emph{et al.}, ``Reward shaping in DRL: A novel framework for adaptive resource management in dynamic environments,'' \emph{Inf. Sci. (N. Y).}, vol. 715, Oct. 2025, doi: 10.1016/j.ins.2025.122238.

\bibitem{b19} A. Abderrahmane, H. Drid, and A. Behaz, ``A Survey of Controller Placement Problem in SDN-IoT Network,'' Dec. 01, 2024, \emph{Springer Science and Business Media B.V.} doi: 10.1007/s44227-024-00035-y.

\bibitem{b20} R. Botez, J. Costa-Requena, I. A. Ivanciu, V. Strautiu, and V. Dobrota, ``Sdn-based network slicing mechanism for a scalable 4g/5g core network: A kubernetes approach,'' \emph{Sensors}, vol. 21, no. 11, Jun. 2021, doi: 10.3390/s21113773.

\bibitem{b21} J. M. Fernandez, I. Vidal, and F. Valera, ``Enabling the orchestration of IoT slices through edge and cloud microservice platforms,'' \emph{Sensors (Switzerland)}, vol. 19, no. 13, Jul. 2019, doi: 10.3390/s19132980.

\bibitem{b22} Y. Zhang, C. Li, N. Chen, and P. Zhang, ``Intelligent Requests Orchestration for Microservice Management Based on Blockchain in Software Defined Networking: A Security Guarantee,'' in \emph{2022 IEEE International Conference on Communications Workshops, ICC Workshops 2022}, Institute of Electrical and Electronics Engineers Inc., 2022, pp. 254--259. doi: 10.1109/ICCWorkshops53468.2022.9814536.

\bibitem{b23} F. Tonini, C. Natalino, M. Furdek, C. Raffaelli, and P. Monti, ``Network Slicing Automation: Challenges and Benefits,'' in \emph{2020 24th International Conference on Optical Network Design and Modeling, ONDM 2020}, 2020. doi: 10.23919/ONDM48393.2020.9133004.

\bibitem{b24} S. Sriharsha Rachakonda and R. Lakshmikanth, ``Enhancing Network Security With Qos-Aware Optimization Models In 6G Network Slicing Architectures.'' {[}Online{]}. Available: https://www.theaspd.com/ijes.php

\bibitem{b25} P. Kochovski, U. Paščinski, V. Stankovski, and M. Ciglarič, ``Pareto-Optimised Fog Storage Services with Novel Service-Level Agreement Specification,'' \emph{Applied Sciences (Switzerland)}, vol. 12, no. 7, Apr. 2022, doi: 10.3390/app12073308.

\bibitem{b26} K. Zia, A. Chiumento, P. Havinga, R. Riggio, and Y. Huang, ``QoS Aware Slice Resource Management Using Deep Reinforcement Learning in IoT Networks,'' in \emph{2023 19th International Conference on Distributed Computing in Smart Systems and the Internet of Things (DCOSS-IoT)}, Jan. 2023, pp. 150--154. doi: 10.1109/DCOSS-IoT58021.2023.00035.

\bibitem{b27} S. Wijethilaka and M. Liyanage, ``Survey on Network Slicing for Internet of Things Realization in 5G Networks,'' \emph{IEEE Communications Surveys and Tutorials}, vol. 23, no. 2, pp. 957 -- 994, 2021, doi: 10.1109/COMST.2021.3067807.

\bibitem{b28} F. Mason, G. Nencioni, and A. Zanella, ``Using Distributed Reinforcement Learning for Resource Orchestration in a Network Slicing Scenario,'' \emph{IEEE/ACM Transactions on Networking}, vol. 31, no. 1, pp. 88 -- 102, 2023, doi: 10.1109/TNET.2022.3187310.

\bibitem{b29} S. Venkatapathy, T. Srinivasan, O. S. Lee, R. Jayaraman, H. G. Jo, and I. H. Ra, ``Slice-aware 5G network orchestration framework based on dual-slice isolation and management strategy (D-SIMS),'' \emph{Sci. Rep.}, vol. 14, no. 1, Dec. 2024, doi: 10.1038/s41598-024-68892-9.

\bibitem{b30} Y. Xu and Q. Yu, ``Deep reinforcement learning based computation offloading and resource allocation strategy for maritime internet of things,'' \emph{Computer Networks}, vol. 264, 2025, doi: 10.1016/j.comnet.2025.111221.

\bibitem{b31} H. She, L. Yan, X. An, C. Mao, and Y. Guo, ``Flow persona: A QoS Flow Rule Scheme based on Deep Learning in SDN-IoT,'' \emph{Computer Networks}, vol. 273, 2025, doi: 10.1016/j.comnet.2025.111784.

\bibitem{b32} A. Holscher, M. Asplund, and F. Boeira, ``Evaluation of an SDN-based Microservice Architecture,'' in \emph{Proceedings of the 2022 IEEE International Conference on Network Softwarization: Network Softwarization Coming of Age: New Challenges and Opportunities, NetSoft 2022}, 2022, pp. 151 -- 156. doi: 10.1109/NetSoft54395.2022.9844094.

\bibitem{b33} V. Sittakul, I. Ioannou, P. Nagaradjane, and V. Vassiliou, ``Intelligent congestion control in 5G URLLC Software-Defined Networks using adaptive resource management via Reinforced Dueling Deep Q-Networks,'' \emph{Journal of Network and Computer Applications}, vol. 242, 2025, doi: 10.1016/j.jnca.2025.104276.

\bibitem{b34} M. Zambianco and G. Verticale, ``Intelligent multi-branch allocation of spectrum slices for inter-numerology interference minimization,'' \emph{Computer Networks}, vol. 196, 2021, doi: 10.1016/j.comnet.2021.108254.

\bibitem{b35} M. F. Ul Islam, S. Hossain, M. G. R. Alam, N. Mansoor, and A. Chakrabarty, ``Optimizing network bandwidth slicing identification: NADAM-enhanced CNN and VAE data preprocessing for enhanced interpretability,'' \emph{PLoS One}, vol. 20, no. 10 October, Oct. 2025, doi: 10.1371/journal.pone.0333286.

\bibitem{b36} L. Zhao, Y. Liu, A. Hawbani, N. Lin, W. Zhao, and K. Yu, ``QoS-Aware Multihop Task Offloading in Satellite--Terrestrial Edge Networks,'' \emph{IEEE Internet Things J.}, vol. 11, no. 19, pp. 31453--31466, Oct. 2024, doi: 10.1109/JIOT.2024.3417456.

\bibitem{b37} C. Ssengonzi, O. P. Kogeda, and T. O. Olwal, ``A survey of deep reinforcement learning application in 5G and beyond network slicing and virtualization,'' \emph{Array}, vol. 14, 2022, doi: 10.1016/j.array.2022.100142.

\bibitem{b38} Y. Zhang, C. Li, N. Chen, and P. Zhang, ``Intelligent Requests Orchestration for Microservice Management Based on Blockchain in Software Defined Networking: A Security Guarantee,'' in \emph{2022 IEEE International Conference on Communications Workshops, ICC Workshops 2022}, 2022, pp. 254 -- 259. doi: 10.1109/ICCWorkshops53468.2022.9814536.

\bibitem{b39} T. Mai, H. Yao, N. Zhang, W. He, D. Guo, and M. Guizani, ``Transfer Reinforcement Learning aided Distributed Network Slicing Optimization in Industrial IoT.''

\bibitem{b40} D. H. Abdulazeez and S. K. Askar, ``Offloading Mechanisms Based on Reinforcement Learning and Deep Learning Algorithms in the Fog Computing Environment,'' \emph{IEEE Access}, vol. 11, pp. 12554--12585, 2023, doi: 10.1109/ACCESS.2023.3241881.

\bibitem{b41} C. Tan, ``Comparative Study of Reinforcement Learning Performance Based on PPO and DQN Algorithms,'' \emph{Applied and Computational Engineering}, vol. 175, pp. 30--36, 2025.

~